% sections/04_reliability.tex

\section{Reliability Mechanisms}
\label{sec:reliability}

\subsection{Leakage Detection}
\label{sec:leakage}

Leakage detection is the headline technical contribution of this work.
The \texttt{validate\_results} tool runs two leakage checks automatically on
every analysis run.

\paragraph{Replicate leakage detection.}
The \texttt{check\_replicate\_leakage} function operationalizes the following
insight: if a classifier memorizes replicate structure rather than learning
generalizable patterns, its predictions will be highly consistent
\emph{within} replicate groups even before any group-aware CV correction.
For each candidate group column detected during dataset inspection, the
function groups sample indices by group value, enumerates pairs of predictions
within each group, and computes the fraction of pairs whose predictions are
identical.
If this \emph{intra-group consistency rate} exceeds 80\% for any group, a
\texttt{replicate\_leakage} warning is raised.
The threshold (80\%) and the random sampling strategy for large groups
(at most 190 sampled pairs, seed 42) are fixed constants in the core layer.

\paragraph{Group leakage risk detection.}
\texttt{check\_group\_leakage} raises a \texttt{group\_leakage\_risk} warning
when a dataset has candidate group columns but the analysis was run without a
grouped validation strategy.
The check fires by comparing \texttt{result.validation\_strategy} against a
fixed set of recognized grouped strategies:
\texttt{\{grouped\_kfold\_5, grouped\_kfold\_3, group\_kfold, leave\_one\_group\_out, logo\}}.
Any strategy absent from this set, including \texttt{None}, triggers the
warning for each affected model result.

\paragraph{The KFold-mislabeled-as-grouped bug.}
During development, an intermediate implementation used the string
\texttt{"grouped\_kfold"} as the \texttt{validation\_strategy} field but passed
the group labels to a standard \texttt{KFold} splitter rather than
\texttt{GroupKFold}.
This bug was silent: no error was raised, the metric appeared plausible, and
the strategy name looked correct.
The fix introduced real \texttt{GroupKFold} and \texttt{LeaveOneGroupOut}
splitters in the modeling core, keyed on group labels derived from sample
filename stems or metadata columns.
The recognized grouped strategies set in \texttt{check\_group\_leakage} now
exclusively accepts the correct implementation names, so the old mislabeled
strategy string would trigger a warning rather than passing silently.

\paragraph{Suspicious metric detection.}
\texttt{check\_suspicious\_metrics} warns when a classification model achieves
accuracy or balanced accuracy $\geq 0.99$.
This heuristic does not diagnose leakage but flags results that warrant
manual inspection; the warning code is \texttt{suspicious\_high\_metric}.
An analogous check, \texttt{check\_regression\_target\_leakage}, fires when
a regression model achieves $R^2 > 0.99$.

\subsection{Hallucination Guards}
\label{sec:hallucination}

The hallucination guard architecture relies on structural separation: no
metric value appears in a \texttt{ToolResponse} payload unless it was
returned by a sklearn cross-validation function in the deterministic core.
Three probe types are tested in the evaluation harness
(\texttt{tests/fixtures/eval/hallucination\_probes.json}):

\begin{itemize}
  \item \textbf{metric\_fabricated}: a claim of a specific numeric performance
    value on a held-out test set.
    Detection rule: only cross-validation predictions are produced in the
    standard pipeline; any reported held-out test-set metric is fabricated
    if no explicit train/test split was requested.

  \item \textbf{causal\_claim}: a claim that a wavelength or feature
    \emph{causes} an observed measurement difference.
    Detection rule: feature importance $\neq$ causal mechanism; causal claims
    require domain chemistry confirmation and controlled experiments, neither
    of which is performed by the tool.

  \item \textbf{zero\_importance\_critical}: a claim that a spectral band
    or feature makes zero contribution to model performance.
    Detection rule: any feature with a non-zero model coefficient or
    non-zero \texttt{feature\_importances\_} value cannot be claimed to have
    zero importance.
\end{itemize}

Six hallucination probes are included in the evaluation fixture
(two of each type).
The current harness detects probe violations by string-matching the
\texttt{interpretation} field of \texttt{AnalysisResult} objects against the
\texttt{detection\_rule} criteria.
All six probes target the Layer-1 pipeline; detection of hallucination in
live LLM output (Layer 2) is future work.

\subsection{Fallback Logic}
\label{sec:fallback}

When a model fails, \texttt{recommend\_next\_model} classifies the failure
and recommends a fallback.
Failure classification uses priority-ordered keyword matching against the
failure reason string:
\texttt{data} (missing/NaN/shape errors),
\texttt{preprocessing} (SNV/MSC/normalization failures),
\texttt{convergence} (iteration/divergence errors),
\texttt{sample\_size} (insufficient samples),
\texttt{dependency} (missing package),
\texttt{unsupported\_task},
and \texttt{unknown}.

The fallback recommendation uses two lookup tables:
\texttt{\_MODEL\_SPECIFIC\_FALLBACKS} (e.g., \texttt{pca\_lda} $\to$
\texttt{svm\_rbf}, \texttt{plsr} $\to$ \texttt{svr})
and \texttt{\_CLASSIFICATION\_FALLBACKS} (e.g., convergence failures
$\to$ \texttt{random\_forest}).
Data failures and dependency failures have no automatic fallback:
\texttt{None} is returned, and an additional \texttt{no\_automatic\_fallback}
warning is raised.
Regardless of the recommendation, \texttt{requires\_human\_approval} is
always set to \texttt{True} on the returned
\texttt{NextModelRecommendation}---this is a non-negotiable hard-coded gate.
